\subsection{ESG Disclosure Theory in the Context of ABSA: Linking Voluntary/Mandatory Frameworks to the Analytical Challenge of Distinguishing Commitment from Outcome}

This synthesis connects the core theoretical lenses used to justify and critique ESG disclosure with the concrete analytical challenge faced by aspect-based sentiment analysis (ABSA) in ESG disclosures: reliably distinguishing commitments from outcomes in bilingual (Indonesian/English) corporate reporting. The discussion integrates voluntary and mandatory disclosure dynamics, dominant reporting standards (GRI and SASB), stakeholder theory (Freeman, 1984), and legitimacy theory, and then maps these theories to the ABSA task of separating commitment statements from outcome disclosures in sentiment signals.

\subsubsection{Voluntary vs mandatory disclosure frameworks and the regulatory trend toward mandatory ESG reporting}

Core idea and theoretical framing. Voluntary ESG disclosures have historically allowed firms to signal governance intentions and social responsibility without binding regulatory consequences. However, a rising body of regulatory scholarship emphasizes a shift toward mandatory ESG reporting to reduce information asymmetry, strengthen market integrity, and curb opportunistic signaling. This regulatory trend is argued to mitigate greenwashing risks by elevating baseline disclosures and standardizing content across firms and jurisdictions \cite{10.1002/bse.2657,10.1145/3322905.3322909,10.1002/bsd2.244,10.1145/3167132.3167236}. The literature notes that mandatory regimes may still entail concerns about superficial compliance, hence the need for robust analytical tools to verify material disclosures against regulatory mandates and verifiable performance data \cite{10.1002/bsd2.70172,10.3390/jrfm17080326,10.2139/ssrn.5130854}

Link to ABSA and the commitment/outcome distinction. In a mandatory regime, the containment of disclosures increases the salience of evidence-bearing sentences that indicate genuine outcomes rather than aspirational commitments. ABSA must therefore differentiate sentences that express commitment (e.g., “we will reduce emissions”) from sentences that assert actual outcomes (e.g., “emissions reduced by 15\%”) to avoid misinterpreting regulatory mandates as substantive performance. The literature’s emphasis on materiality and signaling under regulatory convergence underscores the necessity for sentence-level ABSA that can map sentiment to verifiable disclosures and regulatory requirements \cite{10.1002/bsd2.70172,10.1108/sampj-04-2023-0227,10.3390/su15065508,10.60079/aaar.v2i1.166,10.1108/medar-11-2021-1486} .

Indonesian and global context. Global trends toward mandating ESG reporting resonate with Indonesian reforms oriented by OJK and IDX, where regulators encourage structured, traceable ESG disclosures in filings and public communications \cite{10.22617/tcs210472-2,10.1108/mf-01-2023-0020,10.1007/s40804-021-00233-z} . This cross-border alignment strengthens the expectation that ABSA systems must be capable of parsing both mandatory disclosures and supplementary voluntary narratives, and of distinguishing regulatory-quoted commitments from measured outcomes within Indonesian disclosures that may include bilingual Indonesian/English content \cite{10.22617/tcs210472-2,10.1108/mf-01-2023-0020,10.1007/s40804-021-00233-z,10.1016/j.inffus.2025.103073,10.1111/eufm.12509}

Implications for ABSA design. A principled ABSA approach should incorporate a regulatory-aware ontology that identifies sentences tied to regulatory obligations, cross-references their materiality, and distinguishes indicative commitments from reported outcomes. This also motivates calibration against external indicators and targets (e.g., net-zero trajectories, GHG reductions) to strengthen the verifiability of reported outcomes.

\subsubsection{GRI and SASB standards as dominant structuring frameworks}

Core idea and theoretical framing. GRI offers a broad, stakeholder-focused view of ESG impacts, while SASB emphasizes financially material, industry-specific disclosures intended to inform investment decisions. The coexistence and partial overlap of these frameworks have been documented, with debates about materiality, interoperability, and the potential benefits of integrating GRI's breadth with SASB's financial-materiality orientation \cite{10.1108/medar-11-2021-1486,10.1007/978-3-319-04723-2_7,10.1002/bse.2657,10.3390/su11041093,10.1111/ablj.12148,10.1108/ijlma-06-2022-0113}

The literature also addresses how hybrid or integrated reporting can bridge financial and non-financial disclosures, influencing what firms report and how investors interpret those reports \cite{10.3390/su15032610,10.46647/ijetms.2023.v07i01.006}

Link to ABSA and the commitment/outcome distinction. The ABSA task benefits from explicit mapping of sentences to standardized indicator sets and to regulatory concept clusters. Under GRI, sentences may express broad commitments or qualitative narratives about impacts without precise performance data; under SASB, sentences are more likely to anchor in financially material topics and observable outcomes. An ABSA system must therefore learn to disambiguate sentences that reflect aspirational commitments within a GRI-aligned narrative from sentences that report quantifiable outcomes aligned with SASB materiality. This necessitates cross-framework alignment within the ABSA taxonomy and careful sentence-level labeling that preserves the intended regulatory and investor semantics \cite{10.1002/bsd2.70172,10.1108/medar-11-2021-1486,10.1007/978-3-319-04723-2_7,10.3390/su11041093,10.1108/ijlma-06-2022-0113} 

Indonesian context. Indonesia’s disclosures under OJK/IDX often reference multiple standards and may harmonize GRI and SASB concepts in a single document. An ABSA pipeline operating on bilingual Indonesian/English reports must accommodate cross-language terminologies and indicator mappings, ensuring that sentiment signals attach to the correct framework concept and thereby support regulatory compliance and cross-market comparability \cite{10.22617/tcs210472-2,10.1108/mf-01-2023-0020,10.1007/s40804-021-00233-z,10.1016/j.inffus.2025.103073,10.1111/eufm.12509}

Implications for ABSA design. The ABSA system should facility a dual-maxis taxonomy that (i) anchors sentences to GRI-based materiality and stakeholder signals, and (ii) ties financially material SASB-like indicators to outcome statements and KPIs. A transparent mapping between aspects, signals, and regulatory anchors will improve interpretability and auditability, addressing both signaling dynamics and the regulatory enforcement dimension.

\subsubsection{Stakeholder theory as a rationale for ESG disclosure (Freeman, 1984)}

Core idea and theoretical framing. Freeman’s stakeholder theory posits that firms ought to consider the interests of a broad set of stakeholders beyond shareholders, including employees, customers, suppliers, regulators, communities, and the environment. This theory provides the normative justification for ESG disclosure as a mechanism to manage stakeholder expectations, legitimacy, and social license to operate. In empirical terms, it has been used to explain why firms publish ESG disclosures to signal responsiveness, reduce information asymmetry, and facilitate stakeholder engagement, thereby potentially stabilizing firm performance and reputation \cite{10.3390/jrfm16100429,10.1002/csr.70116,10.1108/sampj-10-2020-0358,10.2139/ssrn.5130854}
Link to ABSA and the commitment/outcome distinction. From a ABSA perspective, commitment statements can reflect stakeholders’ demands and management’s stated intentions to address stakeholder concerns, while outcomes reveal actual progress toward meeting those stakeholder expectations. Distinguishing commitment from outcome is thus central to measuring genuine stakeholder responsiveness versus performative signaling. ABSA that can separate stakeholder-centric commitments (e.g., “we will improve labor practices”) from actual outcomes (e.g., “occupational injury rates fell by X%”) provides a rigorous instrument to assess whether disclosures satisfy stakeholder expectations or primarily signal intent.
Implications for ABSA design. Embedding Freeman’s theory into ABSA encourages explicit annotation of stakeholder focus areas and their corresponding signals, enabling analyses of whether disclosures reflect stakeholder-validated progress rather than generic or rhetorical statements. This alignment supports stakeholder governance research and helps practitioners interpret ESG narratives in terms of accountability to diverse constituencies.
\subsubsection{Legitimacy theory explaining why companies report favorably on ESG topics}

Core idea and theoretical framing. Legitimacy theory argues that firms disclose information to gain, maintain, or restore legitimacy in the eyes of audiences that constrain or enable ongoing operations, often by signaling alignment with social norms, regulatory expectations, and market ideals. This logic has been used to explain the prevalence of ESG reporting and the tendency to present favorable narratives, especially when legitimacy pressures are high or when reporting standards are evolving. Critics note potential weaknesses, such as selective disclosure and performative narratives that do not reflect actual performance,,.
Link to ABSA and the commitment/outcome distinction. Legitimacy-driven disclosures tend to emphasize positive commitments and aspirational statements designed to secure continued social approval. ABSA can be used to interrogate these narratives by analyzing whether commitment claims are substantiated by subsequent outcomes and performance indicators within the text, thereby detecting potential legitimacy-seeking signals that are not corroborated by concrete results. Consequently, sentence-level appraisal of commitment statements versus described outcomes is a critical test of legitimacy signaling versus substantive progress.
Implications for ABSA design. A legitimacy-aware ABSA framework should incorporate detection of legitimacy strategies (e.g., broad social narratives, selective disclosure) and provision of evidence-based assessments by linking commitments to evidence of outcomes and independent verifications. This would enable analysts to distinguish legitimate progress from performative rhetoric and to identify potential greenwashing indicators embedded in seemingly positive disclosures.
\subsubsection{Integrating the theory-to-ABSA mapping: distinguishing commitment from outcome}

Analytical challenge. Across these theories, a central analytical challenge is distinguishing commitment signals (forward-looking or aspirational statements) from outcomes (past or current performance). This distinction is essential because commitments may reflect strategic intent or compliance with normative expectations, whereas outcomes provide observable evidence of impact. The ABSA task thus requires sentence-level disambiguation, robust cross-language term mappings, and the ability to tie linguistic cues to quantifiable performance indicators and regulatory criteria.
Proposed integration approach. A theory-informed ABSA architecture would (i) embed a theory-aligned annotation schema that tags sentences by whether they express commitment, action, or outcome, (ii) map these tokens to GRI/SASB indicators and to stakeholder and legitimacy cues, and (iii) incorporate cross-language lexicons and regulatory anchors to ensure consistent interpretation across Indonesian/English texts. The outcome is an ABSA system that not only detects sentiment polarity at the aspect level but also classifies the epistemic modality (commitment vs outcome) that underpins each sentiment expression.
Implications for research and practice. This integrated approach supports more rigorous enforcement of regulatory disclosures, more credible investor signals, and deeper academic inquiry into the dynamics of ESG communication. It aligns with the broader literature on materiality, signaling, and legitimacy while providing a technically operational framework for ABSA in multilingual ESG reporting.
\subsubsection{Concluding synthesis: theoretical grounding for ABSA in Indonesian ESG disclosures}

The convergence of voluntary/mandatory disclosure trends, dominant frameworks, and foundational theories explains why ABSA must move beyond global polarity to deliver sentence- and topic-level evidence about commitments and outcomes. By anchoring analysis in GRI/SASB, and leveraging stakeholder and legitimacy theories to interpret signaling choices, ABSA can provide granular, auditable insights into whether ESG disclosures reflect substantive progress or rhetorical positioning. The bilingual Indonesian/English reporting context intensifies the need for cross-language alignment and regulatory mapping, ensuring that commitment/outcome distinctions hold across languages and regulatory regimes. This theoretical synthesis provides a rigorous foundation for the methodological choices in the ABSA pipeline and informs the design of annotation schemes, models, and evaluation strategies.